%%%%%%%%%%%%%%%%%%%%%%%%%%%%%%%%%%%%

\section{負の二項分布}

%%%%%%%%%%%%%%%%%%%%%%%%%%%%%%%%%%%%

\begin{frame}
\frametitle{負の二項分布}

\begin{itemize}

\item \hl{負の二項分布}は、$n$ 回目の試行で $k$ 回目の成功が起こる確率を記述する。

\item 負の二項分布の状況を識別するのに役立つ4つの条件：
\begin{enumerate}
\item 試行は独立である。
\item 各試行の結果は成功または失敗に分類される。
\item 成功確率（$p$）は各試行で同じである。
\item 最後の試行は成功でなければならない。
\end{enumerate}
最初の3つの条件は二項分布と共通している。

\end{itemize}

\vfill

\formula{負の二項分布}
{
P($n$ 回目の試行で $k$ 回目の成功) = ${n-1 \choose k-1}~p^k~(1-p)^{n-k}$、\\
ここで $p$ は個々の試行が成功する確率である。すべての試行は独立と仮定する。
}

\end{frame}

%%%%%%%%%%%%%%%%%%%%%%%%%%%%%%%%%%%%

\begin{frame}[shrink=10]

\dq{心理学研究室でアルバイトをしている大学生が、研究に参加するカップルを10組募集するよう頼まれた。彼女は学生センターの前に立ち、建物を出る5人おきに1人に対して、交際中かどうか、もし交際中であれば重要な他者と一緒に研究に参加したいかを尋ねることにした。このような人を見つける確率が10\%だとする。目標を達成するまでに30人に尋ねる必要がある確率は何か？}

与えられた情報：$p = 0.10$、$k = 10$、$n = 30$。$30$ 回目の試行で $10$ 回目の成功の確率を求めるため、負の二項分布を使用する。

\begin{eqnarray*}
P(\text{30回目で10回目の成功}) &=& {29 \choose 9} \times 0.10^{10} \times 0.90^{20} \\
&=& 10,015,005 \times 0.10^{10} \times 0.90^{20} \\
&=& 0.00012
\end{eqnarray*}

\end{frame}

%%%%%%%%%%%%%%%%%%%%%%%%%%%%%%%%%%%%

\begin{frame}
\frametitle{二項分布 vs. 負の二項分布}

\dq{負の二項分布は二項分布とどのように異なるか？}

\begin{itemize}

\item 二項分布の場合、通常は固定された試行回数があり、代わりに成功の回数を考える。

\item 負の二項分布の場合、固定された回数の成功を観測するまでに何回の試行が必要かを調べ、最後の観測が成功であることが必要である。

\end{itemize}

\end{frame}

%%%%%%%%%%%%%%%%%%%%%%%%%%%%%%%%%%%%

\begin{frame}
\frametitle{練習問題}

\pq{次のうち、目的の確率を計算するために負の二項分布を使う状況はどれか？}

\begin{enumerate}[(a)]

\item 5歳の男の子の身長が42インチより高い確率。

\item 10回のソフトボール投球のうち3回成功する確率。

\item ポーカーでストレートフラッシュの手が配られる確率。

\item 最初のヒットの前に8回ミスをする確率。

\solnMult{8回目の試行で3回目のボールに当たる確率。}

\end{enumerate}

\end{frame}
